% --- Archon physics formalization source begin ---
% archon:physics
% archon:covers PhyXMiniProblems/problem_phyx_mini_0285.lean
% archon:source-report reports/phyx_mini/problem_phyx_mini_0285.source.json

\chapter{Physics problem phyx\_mini\_0285}
\label{ch:PhyXMiniProblems_problem_phyx_mini_0285}

\paragraph{Problem source.}
\#\# Physical scenario

During sporting events within large, densely packed stadiums, spectators will send a wave (or pulse) around the stadium. As the wave reaches a group of spectators, they stand with a cheer and then sit. At any instant, the width \$w\$ of the wave is the distance from the leading edge (people are just about to stand) to the trailing edge (people have just sat down). Suppose a human wave travels a distance of \$853\$ seats around a stadium in \$39\$ s, with spectators requiring about \$1.8\$ s to respond to the wave's passage by standing and then sitting.

\#\# Question

What is the width \$w\$ (in number of seats)?

\#\# Answer choices

- A: A: 37

- B: B: 38

- C: C: 40

- D: D: 39

\#\# Figure caption (auxiliary; use the image as primary evidence)

The image shows a group of six people standing behind a brick wall. They are positioned in a line, and most are raising their hands in a cheering or celebratory gesture. The wall is illustrated with a brick pattern.

Above the people, there is an arrow labeled "v" pointing to the right, suggesting motion or direction. Below the wall, another arrow labeled "w" indicates the width of the wall. Both arrows have their labels oriented in the same direction as the arrows themselves. The people are wearing various colored shirts, providing some contrast against the wall.

\#\# Dataset metadata

Wave Properties; Physical Model Grounding Reasoning, Multi-Formula Reasoning

\paragraph{Recorded answer/context.}
D: D: 39

\paragraph{Figure/image path.}
/root/proposal\_for\_physic/science-mango/phyx\_mini\_run/phyx\_data/test\_image/285.png

\paragraph{Formalization target.}
create a compiling Lean file with sorry bodies at `PhyXMiniProblems/problem\_phyx\_mini\_0285.lean`. The Lean declarations must preserve the physical quantities, dimensions or dimensional roles, figure labels, governing-law hypotheses, and final relation expressed by this problem.
Use Mathlib/Physlib names found through LeanExplore where available. If a physics API is missing, introduce faithful local abstractions rather than scalar placeholder aliases.

\begin{theorem}[Physics formalization target]
\label{thm:physics:phyx_mini_0285:target}
\lean{PhyXMiniProblems.ProblemPhyXMini0285.rounded_width_is_thirty_nine_seats}
\uses{def:physics:phyx-mini-0285:phyxminiproblems-problemphyxmini0285-stadiumhumanwave, def:physics:phyx-mini-0285:phyxminiproblems-problemphyxmini0285-hashumanwavephysics, def:physics:phyx-mini-0285:phyxminiproblems-problemphyxmini0285-matchesproblemdata, def:physics:phyx-mini-0285:phyxminiproblems-problemphyxmini0285-matchessourcefigure}
The assigned autoformalize agent should translate this physics problem into Lean declarations in the covered file, with theorem and lemma proof bodies written as `by sorry`.
\end{theorem}
\begin{proof}
This is an autoformalization task, not a proof task. Produce statements that can later be proved by the physics prover without weakening the physical model.
\end{proof}

% --- Archon declaration topology begin ---
\section{Declaration topology}
\begin{definition}[Length Quantity]
  \label{def:physics:phyx-mini-0285:phyxminiproblems-problemphyxmini0285-lengthquantity}
  \lean{PhyXMiniProblems.ProblemPhyXMini0285.LengthQuantity}
  A stadium distance or one-dimensional edge coordinate
\end{definition}
\begin{proof}
  This is the declared model definition of Length Quantity.
\end{proof}
\begin{definition}[Time Quantity]
  \label{def:physics:phyx-mini-0285:phyxminiproblems-problemphyxmini0285-timequantity}
  \lean{PhyXMiniProblems.ProblemPhyXMini0285.TimeQuantity}
  A duration, whose scalar readout is measured in seconds
\end{definition}
\begin{proof}
  This is the declared model definition of Time Quantity.
\end{proof}
\begin{definition}[Speed Quantity]
  \label{def:physics:phyx-mini-0285:phyxminiproblems-problemphyxmini0285-speedquantity}
  \lean{PhyXMiniProblems.ProblemPhyXMini0285.SpeedQuantity}
  A propagation speed, whose scalar readout is measured in seat spacings per second
\end{definition}
\begin{proof}
  This is the declared model definition of Speed Quantity.
\end{proof}
\begin{definition}[Spectator Boundary State]
  \label{def:physics:phyx-mini-0285:phyxminiproblems-problemphyxmini0285-spectatorboundarystate}
  \lean{PhyXMiniProblems.ProblemPhyXMini0285.SpectatorBoundaryState}
  The state of spectators at either boundary of the moving pulse
\end{definition}
\begin{proof}
  This is the declared model definition of Spectator Boundary State.
\end{proof}
\begin{definition}[Propagation Direction]
  \label{def:physics:phyx-mini-0285:phyxminiproblems-problemphyxmini0285-propagationdirection}
  \lean{PhyXMiniProblems.ProblemPhyXMini0285.PropagationDirection}
  The propagation direction shown by the arrow labelled v
\end{definition}
\begin{proof}
  This is the declared model definition of Propagation Direction.
\end{proof}
\begin{definition}[Figure Label]
  \label{def:physics:phyx-mini-0285:phyxminiproblems-problemphyxmini0285-figurelabel}
  \lean{PhyXMiniProblems.ProblemPhyXMini0285.FigureLabel}
  Symbols printed on the two arrows in the source figure
\end{definition}
\begin{proof}
  This is the declared model definition of Figure Label.
\end{proof}
\begin{definition}[Figure Feature]
  \label{def:physics:phyx-mini-0285:phyxminiproblems-problemphyxmini0285-figurefeature}
  \lean{PhyXMiniProblems.ProblemPhyXMini0285.FigureFeature}
  The two annotated features in the source figure
\end{definition}
\begin{proof}
  This is the declared model definition of Figure Feature.
\end{proof}
\begin{definition}[Width Marker Kind]
  \label{def:physics:phyx-mini-0285:phyxminiproblems-problemphyxmini0285-widthmarkerkind}
  \lean{PhyXMiniProblems.ProblemPhyXMini0285.WidthMarkerKind}
  The width annotation is a two-ended span rather than a direction of travel
\end{definition}
\begin{proof}
  This is the declared model definition of Width Marker Kind.
\end{proof}
\begin{definition}[Source Figure]
  \label{def:physics:phyx-mini-0285:phyxminiproblems-problemphyxmini0285-sourcefigure}
  \lean{PhyXMiniProblems.ProblemPhyXMini0285.SourceFigure}
  \uses{def:physics:phyx-mini-0285:phyxminiproblems-problemphyxmini0285-lengthquantity, def:physics:phyx-mini-0285:phyxminiproblems-problemphyxmini0285-propagationdirection, def:physics:phyx-mini-0285:phyxminiproblems-problemphyxmini0285-figurelabel, def:physics:phyx-mini-0285:phyxminiproblems-problemphyxmini0285-figurefeature, def:physics:phyx-mini-0285:phyxminiproblems-problemphyxmini0285-widthmarkerkind}
  A structured transcription of the primary image
\end{definition}
\begin{proof}
  This is the declared model definition of Source Figure.
\end{proof}
\begin{definition}[Stadium Human Wave]
  \label{def:physics:phyx-mini-0285:phyxminiproblems-problemphyxmini0285-stadiumhumanwave}
  \lean{PhyXMiniProblems.ProblemPhyXMini0285.StadiumHumanWave}
  \uses{def:physics:phyx-mini-0285:phyxminiproblems-problemphyxmini0285-lengthquantity, def:physics:phyx-mini-0285:phyxminiproblems-problemphyxmini0285-timequantity, def:physics:phyx-mini-0285:phyxminiproblems-problemphyxmini0285-speedquantity, def:physics:phyx-mini-0285:phyxminiproblems-problemphyxmini0285-spectatorboundarystate, def:physics:phyx-mini-0285:phyxminiproblems-problemphyxmini0285-propagationdirection, def:physics:phyx-mini-0285:phyxminiproblems-problemphyxmini0285-sourcefigure}
  The physical quantities of the stadium pulse. waveWidth is independent data in the structure: it is not defined as the requested numerical answer. Its relationship to the edges and to the response duration is supplied separately by the governing-law predicate
\end{definition}
\begin{proof}
  This is the declared model definition of Stadium Human Wave.
\end{proof}
\begin{definition}[Satisfies Speed Distance Time Law]
  \label{def:physics:phyx-mini-0285:phyxminiproblems-problemphyxmini0285-satisfiesspeeddistancetimelaw}
  \lean{PhyXMiniProblems.ProblemPhyXMini0285.SatisfiesSpeedDistanceTimeLaw}
  \uses{def:physics:phyx-mini-0285:phyxminiproblems-problemphyxmini0285-stadiumhumanwave}
  The dimensionally typed constant-speed relation v = d / t
\end{definition}
\begin{proof}
  This is the declared model definition of Satisfies Speed Distance Time Law.
\end{proof}
\begin{definition}[Satisfies Width Response Law]
  \label{def:physics:phyx-mini-0285:phyxminiproblems-problemphyxmini0285-satisfieswidthresponselaw}
  \lean{PhyXMiniProblems.ProblemPhyXMini0285.SatisfiesWidthResponseLaw}
  \uses{def:physics:phyx-mini-0285:phyxminiproblems-problemphyxmini0285-stadiumhumanwave}
  During the response interval, the pulse advances by one pulse width: w = v Δt
\end{definition}
\begin{proof}
  This is the declared model definition of Satisfies Width Response Law.
\end{proof}
\begin{definition}[Satisfies Edge Separation Law]
  \label{def:physics:phyx-mini-0285:phyxminiproblems-problemphyxmini0285-satisfiesedgeseparationlaw}
  \lean{PhyXMiniProblems.ProblemPhyXMini0285.SatisfiesEdgeSeparationLaw}
  \uses{def:physics:phyx-mini-0285:phyxminiproblems-problemphyxmini0285-stadiumhumanwave}
  The width w is the separation from the trailing edge to the leading edge
\end{definition}
\begin{proof}
  This is the declared model definition of Satisfies Edge Separation Law.
\end{proof}
\begin{definition}[Has Human Wave Physics]
  \label{def:physics:phyx-mini-0285:phyxminiproblems-problemphyxmini0285-hashumanwavephysics}
  \lean{PhyXMiniProblems.ProblemPhyXMini0285.HasHumanWavePhysics}
  \uses{def:physics:phyx-mini-0285:phyxminiproblems-problemphyxmini0285-stadiumhumanwave, def:physics:phyx-mini-0285:phyxminiproblems-problemphyxmini0285-satisfiesspeeddistancetimelaw, def:physics:phyx-mini-0285:phyxminiproblems-problemphyxmini0285-satisfieswidthresponselaw, def:physics:phyx-mini-0285:phyxminiproblems-problemphyxmini0285-satisfiesedgeseparationlaw}
  The physical laws used to infer the pulse width; no numerical answer occurs here
\end{definition}
\begin{proof}
  This is the declared model definition of Has Human Wave Physics.
\end{proof}
\begin{definition}[Matches Problem Data]
  \label{def:physics:phyx-mini-0285:phyxminiproblems-problemphyxmini0285-matchesproblemdata}
  \lean{PhyXMiniProblems.ProblemPhyXMini0285.MatchesProblemData}
  \uses{def:physics:phyx-mini-0285:phyxminiproblems-problemphyxmini0285-stadiumhumanwave}
  The numerical measurements and boundary meanings stated in the problem. The exact rational 9 / 5 is the 1.8 s response-time readout
\end{definition}
\begin{proof}
  This is the declared model definition of Matches Problem Data.
\end{proof}
\begin{definition}[Matches Source Figure]
  \label{def:physics:phyx-mini-0285:phyxminiproblems-problemphyxmini0285-matchessourcefigure}
  \lean{PhyXMiniProblems.ProblemPhyXMini0285.MatchesSourceFigure}
  \uses{def:physics:phyx-mini-0285:phyxminiproblems-problemphyxmini0285-stadiumhumanwave}
  The figure shows six spectators, a rightward arrow labelled v, and a width span labelled w whose endpoints coincide with the pulse edges
\end{definition}
\begin{proof}
  This is the declared model definition of Matches Source Figure.
\end{proof}
% --- Archon declaration topology end ---
% --- Archon physics formalization source end ---
